\section{Introduction}
UAV swarms require frequent exchange of state, control, and coordination data, often outside reliable infrastructure coverage. NR sidelink is therefore a natural candidate for direct intra-swarm communication \cite{mishra2022cu2x,varonen2024sidelink}. Aerial line-of-sight propagation is beneficial for desired links, but it can also expose a receiver to many strong simultaneous interferers. Prior work has studied multi-UAV outage \cite{lau2023outage}, joint trajectory/frequency/routing optimization \cite{alam2024jtfr}, site-specific unmanned-platform sidelink \cite{giannakoulas2025uxu}, and synchronization/access design \cite{zhang2026sync}. What remains less explicit is how swarm topology interacts with the bandwidth cost of resource isolation and NR-aware link reliability.

This paper asks: \emph{How do swarm density and communication topology change the NR sidelink interference regime, and how far can resource separation and directional spatial selectivity extend the feasible reliability/goodput region?} The contributions are threefold. First, we combine measurement-derived A2A propagation with standards-based NR transport-block mechanics and sourced link-level BLER curves. Second, we quantify an exact PRB-partition operating envelope under a controlled full-load baseline. Third, we test whether the headline conclusions survive changes in pairing, resource allocator, link adaptation, and noise-bandwidth accounting. This last step is essential: the severe density collapse observed under arbitrary long desired links is not topology invariant.
